\subsection{GRASSMANNIANS}

If K is a field and X is any matrix of \(\mathrm{L}_{n+1,p+1}(\mathrm{K})\) , let \(c_{\sigma}(X)\) denote the determinant of \(X_{\sigma}\) ; in this way, to each matrix X of \(\mathrm{L}_{n+1,p+1}(\mathrm{K})\) correspond

\[
h = \binom {n + 1} {p + 1}
\]

determinants, not all zero (the components of the exterior product of the \(p + 1\) rows of \(X\) ). If we make correspond to \(X\) the point of the projective space \(\mathbf{P}_{h-1}(\mathbf{K})\) whose homogeneous coordinates are the \(c_{\sigma}(X)\) , we have defined a mapping of \(\mathbf{L}_{n+1,p+1}(\mathbf{K})\) into \(\mathbf{P}_{h-1}(\mathbf{K})\) , compatible with the relation \(\Delta_{n,p}(\mathbf{K})\) ; passing to the quotient, we have therefore a mapping \(f\) of \(\mathbf{P}_{n,p}(\mathbf{K})\) into \(\mathbf{P}_{h-1}(\mathbf{K})\) . The image \(\mathbf{G}_{n,p}(\mathbf{K})\) of \(\mathbf{P}_{n,p}(\mathbf{K})\) under this mapping is called the Grassmannian of indices \(n\) , \(p\) . We recall also that the mapping \(f\) is injective, for if \(X\) is a matrix such that \(X_{\sigma}\) is non-singular, the matrix \(\Upsilon = X_{\sigma}^{-1}X\) of \(\mathbf{B}_{\sigma}\) which corresponds to the class of \(X\) mod. \(\Delta_{n,p}(\mathbf{K})\) is the matrix \((d_{ij}/c_{\sigma}(X))\) ( \(1 \leqslant i \leqslant p + 1\) , \(0 \leqslant j \leqslant n\) ), where \(d_{ij}\) denotes the determinant of the matrix obtained from \(X_{\sigma}\) by replacing the ith column of \(X_{\sigma}\) by the jth column of X {[}which implies that \(d_{ij}\) is, up to sign, equal to one of the \(c_{\tau}(X)\) {]}.

When \(\mathbf{K}\) is the field \(\mathbf{R}\) , this mapping \(f\) is evidently continuous. The inverse mapping \(g\) is also continuous; for the elements of a matrix belonging to \(\mathbf{B}_{\sigma}\) are rational functions of the homogeneous coordinates of the point of the Grassmannian to which it corresponds; since \(f(\mathbf{B}_{\sigma}) = \mathbf{B}_{\sigma}'\) is the set of points of \(\mathbf{G}_{n,p}\) whose homogeneous coordinate with index \(\sigma\) is not \(0\) , it is an open set in \(\mathbf{G}_{n,p}\) ; hence \(g\) is continuous at every point of \(\mathbf{B}_{\sigma}'\) , and since every point of \(\mathbf{G}_{n,p}\) belongs to at least one set \(\mathbf{B}_{\sigma}'\) , \(g\) is continuous at every point. Thus:

PROPOSITION 9. The Grassmannian \(\mathbf{G}_{n,p}\) is homeomorphic to the space \(\mathbf{P}_{n,p}\) .

We recall finally that the Grassmannians \(\mathbf{G}_{n,p}(\mathbf{K})\) and \(\mathbf{G}_{n,n-p-1}(\mathbf{K})\) are subsets of the same projective space \(\mathbf{P}_{h-1}(\mathbf{K})\) and can be transformed into the other by a projective linear transformation; it follows that \(G_{n,p}\) and \(G_{n,n-p-1}\) are homeomorphic.
% semantic-fragments: legacy-input-seam
\relax{}
\ExerciseGroupsStart
